\section*{Capítulo \ref{ch:b2:heart} --- O coração e o ciclo cardíaco}

\begin{solution}{exo:b2:heart:1}
\omterm{def:b2:blood-circulation:vessels}{Veia} cava $\to$ \omterm{def:b2:heart:anatomy}{átrio} direito $\to$ valva tricúspide $\to$ \omterm{def:b2:heart:anatomy}{ventrículo}
direito $\to$ valva pulmonar $\to$ \omterm{def:b2:blood-circulation:vessels}{artéria} pulmonar $\to$ pulmões
$\to$ \omterm{def:b2:blood-circulation:vessels}{veias} pulmonares $\to$ \omterm{def:b2:heart:anatomy}{átrio} esquerdo $\to$ valva mitral $\to$
\omterm{def:b2:heart:anatomy}{ventrículo} esquerdo $\to$ valva aórtica $\to$ aorta.
\end{solution}

\begin{solution}{exo:b2:heart:2}
Enchimento (diástole): valvas AV abertas, semilunares fechadas.
Contração isovolumétrica: as quatro fechadas. Ejeção: semilunares
abertas, AV fechadas. Relaxamento isovolumétrico: todas fechadas.
Primeira bulha: o fechamento das valvas AV no início da sístole;
segunda: o fechamento das valvas semilunares no fim.
\end{solution}

\begin{solution}{exo:b2:heart:3}
\omterm{prop:b2:heart:conduction}{Nó sinoatrial} ($t = 0$; a onda P, enquanto os \omterm{def:b2:heart:anatomy}{átrios} se despolarizam);
\omterm{prop:b2:heart:conduction}{nó atrioventricular} (alcançado em cerca de \qty{80}{ms}, retendo o
impulso por \qty{100}{ms}: o segmento PR); feixe de His e \omterm{prop:b2:heart:conduction}{fibras de
Purkinje} (\qtyrange{180}{220}{ms}; o QRS, enquanto os \omterm{def:b2:heart:anatomy}{ventrículos} se
despolarizam); a onda T mais tarde, enquanto eles se repolarizam.
\end{solution}

\begin{solution}{exo:b2:heart:4}
Dentro de sua faixa de funcionamento, o \omterm{def:b2:heart:anatomy}{ventrículo} ejeta mais quando
foi mais cheio. Se o \omterm{def:b2:heart:anatomy}{ventrículo} direito bombeia mais, o esquerdo recebe
mais alguns batimentos depois e, pela mesma lei, bombeia mais: todo
desequilíbrio se corrige sem nenhum sinal.
\end{solution}

\begin{solution}{exo:b2:heart:5}
\omterm{prop:b2:heart:cycle}{Volume sistólico} \qty{80}{mL}; \omterm{prop:b2:heart:cycle}{fração de ejeção} $80/130 = \qty{62}{\%}$;
débito $80\times 70 = \qty{5.6}{L/min}$. Em insuficiência:
\qty{60}{mL}, $60/180 = \qty{33}{\%}$, \qty{4.2}{L/min} --- o \omterm{def:b2:heart:anatomy}{ventrículo}
dilatado mantém o débito quase normal trabalhando com um grande volume,
mas ejeta apenas um terço do que contém.
\end{solution}

\begin{solution}{exo:b2:heart:6}
$W = 100\times 133\times 7\times 10^{-5} = \qty{0.93}{J}$; potência
$0.93\times 1.2 = \qty{1.1}{W}$. A \qty{150}{mmHg}: \qty{1.4}{J},
\qty{1.7}{W}. Os \qty{0.56}{W} mecânicos a mais são \qty{3.7}{W}
metabólicos, \qty{0.19}{mL} de oxigênio por segundo: \qty{11}{mL/min} a
mais.
\end{solution}

\begin{solution}{exo:b2:heart:7}
$60/0.5 = 120$ batimentos por minuto; $60/1.5 = 40$. Terceiro: bloqueio
cardíaco completo --- os \omterm{def:b2:heart:anatomy}{átrios} batem a 100 sob o comando do nó, os
\omterm{def:b2:heart:anatomy}{ventrículos} a 33 sob um marca-passo próprio, e o nó AV não conduz nada.
\end{solution}

\begin{solution}{exo:b2:heart:8}
Subida $20/25 = \qty{0.8}{s}$, mais \qty{0.15}{s}: \qty{0.95}{s}, 63
batimentos por minuto. Acetilcolina: $25/18 = \qty{1.39}{s}$, mais
$0.15$: \qty{1.54}{s}, 39 por minuto. Noradrenalina:
$20/40 = \qty{0.5}{s}$, mais $0.15$: \qty{0.65}{s}, 92 por minuto.
\end{solution}

\begin{solution}{exo:b2:heart:9}
O atraso permite que os \omterm{def:b2:heart:anatomy}{átrios} terminem de se esvaziar nos \omterm{def:b2:heart:anatomy}{ventrículos}
antes que os \omterm{def:b2:heart:anatomy}{ventrículos} se contraiam; sem ele, \omterm{def:b2:heart:anatomy}{átrios} e \omterm{def:b2:heart:anatomy}{ventrículos} se
contrairiam juntos, as valvas AV se fechariam sobre \omterm{def:b2:heart:anatomy}{ventrículos} meio
cheios e o \omterm{prop:b2:heart:cycle}{volume sistólico} cairia o equivalente à contribuição atrial.
Um intervalo PR longo apenas adia a sístole ventricular e não custa
nada, até que o atraso se torne tão longo que batimentos falhem ou a
condução cesse por completo.
\end{solution}

\begin{solution}{exo:b2:heart:10}
Repouso: $5000/45 = \qty{111}{mL}$; exercício: $\num{30000}/180 =
\qty{167}{mL}$. A lei de Starling converte o maior retorno venoso do
exercício num \omterm{prop:b2:heart:cycle}{volume sistólico} maior; os nervos simpáticos acrescentam
frequência e contratilidade (um \omterm{prop:b2:heart:cycle}{volume sistólico} final menor). Acima de
cerca de 200, a diástole é curta demais para encher o \omterm{def:b2:heart:anatomy}{ventrículo} e o
\omterm{prop:b2:heart:cycle}{volume sistólico} cai mais depressa do que a frequência sobe.
\end{solution}

\begin{solution}{exo:b2:heart:11}
$v = Q/A$: $250/3 = \qty{83}{cm/s}$ e $250/1 = \qty{250}{cm/s}$.
$\Delta P = \tfrac12\times 1060\times v^{2}$: \qty{365}{Pa}
(\qty{2.7}{mmHg}) e \qty{3300}{Pa} (\qty{25}{mmHg}). O \omterm{def:b2:heart:anatomy}{ventrículo} gera
a pressão extra espessando sua parede; o \omterm{def:b2:heart:anatomy}{ventrículo} espesso e rígido
acaba superando sua irrigação coronária, enche-se mal e entra em
insuficiência.
\end{solution}

\begin{solution}{exo:b2:heart:12}
O automatismo dá um batimento sem nenhum estímulo externo; a lei de
Starling ajusta o débito ao retorno e equilibra os dois \omterm{def:b2:heart:anatomy}{ventrículos} sem
nenhum estímulo externo; os nervos e a adrenalina apenas fixam a
frequência e a contratilidade em torno desse núcleo autorregulado. Um
\omterm{def:b2:heart:anatomy}{coração} transplantado bate (a cerca de 100, sem tônus vagal), ajusta seu
\omterm{prop:b2:heart:cycle}{volume sistólico} ao enchimento e consegue elevar seu débito no
exercício pelo mecanismo de Starling e pela adrenalina circulante ---
mas lentamente, e menos que um \omterm{def:b2:heart:anatomy}{coração} inervado.
\end{solution}

\begin{solution}{pb:b2:heart:1}
\textbf{1.} \qty{70}{mL}; $70/130 = \qty{54}{\%}$; $70\times 70 =
\qty{4.9}{L/min}$.
\textbf{2.} $60/70 = \qty{0.86}{s}$; diástole $0.86 - 0.3 =
\qty{0.56}{s}$.
\textbf{3.} $70/0.3 = \qty{233}{mL/s}$; $233/3 = \qty{78}{cm/s}$.
\textbf{4.} $70/1500 = \qty{0.047}{s}$, cerca de \qty{50}{ms}.
\textbf{5.} $100\times 133\times 7\times 10^{-5} = \qty{0.93}{J}$;
$0.93\times 70/60 = \qty{1.1}{W}$.
\textbf{6.} Os dois \omterm{def:b2:heart:anatomy}{ventrículos} $1.1\times 1.2 = \qty{1.3}{W}$; a
\qty{15}{\%}, \qty{8.7}{W} metabólicos; $8.7/20 = \qty{0.43}{mL}$ de
oxigênio por segundo, \qty{26}{mL/min}.
\textbf{7.} Extraindo \qty{0.14}{mL} por mililitro de \omterm{def:b2:blood-circulation:blood}{sangue}:
$26/0.14 = \qty{190}{mL/min}$.
\textbf{8.} Intervalo de \qty{0.86}{s}, tempo de subida de
\qty{0.71}{s}: $20/0.71 = \qty{28}{mV/s}$.
\textbf{9.} Intervalo de \qty{0.6}{s}, subida de \qty{0.45}{s}:
\qty{44}{mV/s}.
\textbf{10.} Intervalo de \qty{0.33}{s}, subida de \qty{0.18}{s}:
\qty{110}{mV/s}.
\textbf{11.} R--R \qty{0.86}{s}; P: despolarização atrial; QRS:
despolarização ventricular; T: repolarização ventricular.
\textbf{12.} O atraso permite que os \omterm{def:b2:heart:anatomy}{átrios} se esvaziem nos \omterm{def:b2:heart:anatomy}{ventrículos}
antes que estes se contraiam; a 180, o ciclo inteiro dura
\qty{0.33}{s}, de modo que um atraso inalterado de \qty{0.16}{s}
consumiria metade dele --- os nervos simpáticos aceleram também a
condução do nó.
\textbf{13.} Bloqueio cardíaco completo: os \omterm{def:b2:heart:anatomy}{ventrículos} a 40 por
minuto, comandados pelo feixe ou pelas \omterm{prop:b2:heart:conduction}{fibras de Purkinje}, os \omterm{def:b2:heart:anatomy}{átrios} a
75 sob o comando do \omterm{prop:b2:heart:conduction}{nó sinoatrial}.
\textbf{14.} \qty{120}{mL}; $120/150 = \qty{80}{\%}$; $120\times 180 =
\qty{21.6}{L/min}$.
\textbf{15.} $0.33 - 0.3 = \qty{0.03}{s}$: quase nenhum tempo para
encher; mais rápido ainda, e o \omterm{prop:b2:heart:cycle}{volume sistólico} despenca.
\textbf{16.} $W = 120\times 133\times 1.2\times 10^{-4} = \qty{1.9}{J}$,
$\times 3$ por segundo: \qty{5.7}{W}; os dois \omterm{def:b2:heart:anatomy}{ventrículos},
\qty{6.9}{W}; metabólicos, \qty{46}{W}; oxigênio, \qty{2.3}{mL/s},
\qty{140}{mL/min}.
\textbf{17.} $140/0.14 = \qty{1000}{mL/min}$, cinco vezes o fluxo de
repouso, a ser fornecido numa diástole reduzida a um décimo do ciclo:
as \omterm{def:b2:blood-circulation:vessels}{arteríolas} coronárias precisam se dilatar ao máximo.
\textbf{18.} Frequência de cerca de 100 sem tônus vagal, \omterm{prop:b2:heart:cycle}{volume
sistólico} final inalterado em \qty{60}{mL}: \omterm{prop:b2:heart:cycle}{volume sistólico}
$150 - 60 = \qty{90}{mL}$, débito de \qty{9}{L/min}.
\textbf{19.} Só a frequência ($70 \to 180$ a \qty{70}{mL}): $+7.7$;
o enchimento (\qty{20}{mL} a mais a 180): $+3.6$; a contratilidade
(\qty{30}{mL} de resíduo a menos a 180): $+5.4$ --- de 4.9 a
\qty{21.6}{L/min} ($4.9 + 7.7 + 3.6 + 5.4 = 21.6$).
\textbf{20.} $233/1 = \qty{233}{cm/s}$, \qty{2.3}{m/s}.
\textbf{21.} $\tfrac12\times 1060\times 2.33^{2} = \qty{2900}{Pa}$,
\qty{22}{mmHg}.
\textbf{22.} $120/0.3 = \qty{400}{mL/s}$, \qty{4}{m/s}; $\tfrac12\times
1060\times 16 = \qty{8500}{Pa}$, \qty{64}{mmHg}.
\textbf{23.} Pico de $120 + 22 = \qty{142}{mmHg}$ em repouso, cerca de
$140 +
64 = \qty{204}{mmHg}$ no exercício; a pressão média de ejeção em repouso
sobe de 100 para 122: trabalho por batimento $\times 1.22$.
\textbf{24.} Uma parede mais espessa gera mais força e reduz a tensão
por fibra; mas a massa a irrigar cresce, enquanto o fluxo coronário,
comprimido pela parede espessa na sístole e com uma diástole mais curta,
não acompanha, e uma parede espessa relaxa mal e se enche menos:
seguem-se a isquemia e a insuficiência.
\textbf{25.} \qty{4.9}{L/min} em repouso, \qty{21.6}{L/min} no
exercício; \qty{0.93}{J} por batimento; gradiente de \qty{22}{mmHg} em
repouso e de \qty{64}{mmHg} no exercício.
\end{solution}
